% sections/02_science.tex
% 科学层分析：Q1--Q3

\section{科学层：散热物理可行性验证}

科学层的任务是回答一个基础性问题：马斯克设想的太空算力集群，是否符合已知物理定律？我们以斯特藩-玻尔兹曼定律（Stefan-Boltzmann Law）为核心工具，交叉验证AI1卫星的公开数据，并评估将散热规模扩展至GW量级时面临的物理边界。

\subsection{斯特藩-玻尔兹曼定律：T\textsuperscript{4}的非线性约束}

辐射冷却的基本方程为：
\begin{equation}
q = \varepsilon\sigma T^{4}
\end{equation}
其中：$q$——散热通量（\unitWpm{}）；$\varepsilon$——表面发射率（$0<\varepsilon\leq 1$，$\varepsilon=1$为理想黑体）；$\sigma=5.67\times10^{-8}$\,W\,m$^{-2}$\,K$^{-4}$为斯特藩-玻尔兹曼常数；$T$——辐射面温度（K）。

该方程的核心含义：散热能力与温度的四次方成正比。这意味着温度从300\,K提升至400\,K（仅33\%），散热能力翻倍。但代价是GPU结温同步升高，进而导致漏电流增加与器件老化加速。

图~\ref{fig:sb_t4} 绘制了不同发射率下的散热曲线。AI1卫星设计工作点（$T\approx342.5$\,K，$q\approx1400$\,\unitWpm{}，$\varepsilon\approx 0.92$）位于$\varepsilon=0.90$与$\varepsilon=0.95$曲线之间，处于工程可实现范围内，但余量有限。

\begin{figure}[H]
\centering
% figures/sb_t4_curve.tex
% 斯特藩-玻尔兹曼 T⁴ 曲线：散热通量 vs 温度
\begin{tikzpicture}
\begin{axis}[
    width=0.95\textwidth,
    height=7cm,
    xlabel={\textbf{散热面温度} $T$ (K)},
    ylabel={\textbf{散热通量} (W/m\textsuperscript{2})},
    xmin=290, xmax=510,
    ymin=0, ymax=3800,
    grid=major,
    legend style={
        at={(0.02,0.98)},
        anchor=north west,
        font=\footnotesize,
        cells={anchor=west}
    },
    tick label style={font=\footnotesize},
    label style={font=\small},
]

% ε=0.85
\addplot[color=gray!60, line width=1.2pt, mark=*, mark size=1.5pt]
    coordinates {(300,390)(323,523)(340,642)(350,723)(373,932)(400,1234)(450,1981)(500,3015)};
\addlegendentry{$\varepsilon=0.85$}

% ε=0.90
\addplot[color=accent, line width=1.2pt, mark=square*, mark size=1.5pt]
    coordinates {(300,413)(323,555)(340,682)(350,766)(373,988)(400,1306)(450,2093)(500,3190)};
\addlegendentry{$\varepsilon=0.90$}

% ε=0.95
\addplot[color=primary!70, line width=1.2pt, mark=triangle*, mark size=1.5pt]
    coordinates {(300,436)(323,586)(340,720)(350,808)(373,1043)(400,1379)(450,2209)(500,3367)};
\addlegendentry{$\varepsilon=0.95$}

% ε=1.00 (黑体)
\addplot[color=primary, line width=1.5pt, dashed, mark=diamond*, mark size=1.5pt]
    coordinates {(300,459)(323,617)(340,758)(350,851)(373,1097)(400,1452)(450,2325)(500,3544)};
\addlegendentry{$\varepsilon=1.00$ (黑体)}

% AI1 工作点标注
\addplot[only marks, mark=*, mark size=4pt, color=highlight]
    coordinates {(342.5,1400)};
\node[highlight, anchor=south west, font=\footnotesize\bfseries]
    at (axis cs:342.5,1400) {AI1 工作点 (342.5\,K, $\sim$1400\,W/m\textsuperscript{2})};

% T⁴ 非线性辅助标注
\draw[->, >=stealth, graybase, line width=0.6pt]
    (axis cs:380,800) -- (axis cs:460,1800);
\node[graybase, font=\footnotesize, anchor=north]
    at (axis cs:420,1050) {\textbf{$T^{4}$ 非线性增长}};

\end{axis}
\end{tikzpicture}

\caption{斯特藩-玻尔兹曼 T\textsuperscript{4} 散热曲线与AI1工作点}
\label{fig:sb_t4}
\end{figure}

从图~\ref{fig:sb_t4} 中可直观看到T\textsuperscript{4}增长效应：在350\,K附近，温度每升高10\,K，散热通量增加约12\%--13\%；而在450\,K以上，同样的10\,K增量带来的散热增量已接近20\%。这意味着在高温段，散热效率本身不是问题，问题变成了GPU能否承受对应的高结温。

\subsection{AI1卫星1400\,\unitWpm{}散热密度验证}

SpaceX公布的AI1卫星散热密度约为1400\,\unitWpm{}\cite{spacex2026ai1}。我们使用黑体辐射方程$T=(q/\varepsilon\sigma)^{1/4}$反算。

\begin{table}[H]
\centering
\footnotesize
\caption{AI1散热密度（1400\,\unitWpm{}）反算验证}
\label{tab:ai1_verify}
\begin{tabular}{cccc}
\toprule
$\varepsilon$ & $T$ (K) & $T$ (\unitC{}) & 验证结论 \\
\midrule
1.00（黑体）& 397 & 124 & 低于AI1公开温度 $\rightarrow$ 排除纯黑体 \\
0.95          & 402 & 129 & 接近但偏高 \\
\textbf{0.92} & \textbf{405} & \textbf{132} & \textbf{与AI1设计温度一致} $\checkmark$ \\
0.85          & 412 & 139 & 偏高 \\
\bottomrule
\end{tabular}
\end{table}

反算结果（表~\ref{tab:ai1_verify}）显示，1400\,\unitWpm{}对应$\varepsilon\approx 0.92$、$T\approx 405$\,K（132\,\unitC{}），与SpaceX公开的AI1散热设计参数一致，表明该工作点在物理上成立，发射率参数也在当前航天热控涂层技术水平之内\cite{nasa1993thermal}。

\subsection{1\,GW规模散热面积推算}

马斯克提出的1\,GW年部署量\cite{musk2026terafab}意味着什么？我们进行简单的面积换算。

\begin{table}[H]
\centering
\footnotesize
\caption{不同散热密度下1\,GW的散热面积需求}
\label{tab:gw_area}
\begin{tabular}{ccc}
\toprule
散热密度 $q$（\unitWpm{}） & 所需面积（m\textsuperscript{2}） & 可视化参考 \\
\midrule
1400（AI1基准） & $7.14\times 10^{5}$ & 100个标准足球场 \\
2000（高温方案） & $5.00\times 10^{5}$ & 70个标准足球场 \\
4000（尖端材料\cite{kanti2025graphene}） & $2.50\times 10^{5}$ & 35个标准足球场 \\
6000（理论极限） & $1.67\times 10^{5}$ & 23个标准足球场 \\
\bottomrule
\end{tabular}
\end{table}

在AI1基准散热密度下，1\,GW集群需要约71万平方米的辐射散热面——相当于100个标准足球场的面积。即使采用最乐观的尖端石墨烯复合材料方案\cite{kanti2025graphene}（$q\approx4000$\,\unitWpm{}），仍需约25万平方米。这不是材料层面的"能否制造"问题，而是航天器在轨展开、姿态控制与遮挡管理的系统工程挑战。

\subsection{LEO轨道热环境：太空不是"理想冷沉"}

"太空是巨大的冷沉"这一直觉描述在传播中产生了误导。图~\ref{fig:leo_thermal} 展示了低地球轨道的真实热环境。

\begin{figure}[H]
\centering
% figures/leo_thermal_env.tex
% LEO 轨道热环境示意图 —— 纯 TikZ
\begin{tikzpicture}[
    >=Stealth,
    scale=1.0,
    every node/.style={font=\footnotesize},
]

% ===== 中心卫星 =====
\filldraw[fill=primary!10, draw=primary, line width=1.5pt]
    (-1.5,-0.8) rectangle (1.5,0.8);
\node[primary, font=\bfseries] at (0,0) {\textbf{AI1 散热板}};
\node[primary, font=\scriptsize] at (0,-0.5) {(辐射散热至深空 2.7\,K)};

% ===== 太阳直射（上方金黄色，从太阳指向卫星） =====
\draw[->, line width=1.5pt, color=yellow!80!orange]
    (0,3.5) -- (0,0.8);
\node[anchor=south, text=yellow!60!orange, font=\footnotesize\bfseries]
    at (0,3.6) {\textbf{太阳直射}};
\node[anchor=north, text=yellow!70!orange, font=\scriptsize]
    at (0,2.8) {1361\,W/m\textsuperscript{2}};
\node[anchor=north, text=yellow!70!orange, font=\scriptsize]
    at (0,2.3) {吸收$\sim$5\% (knife-edge 定向后)};

% ===== 地球红外辐射（下方蓝色，从地球指向卫星） =====
\draw[->, line width=1.2pt, color=blue!60]
    (-1.2,-3.0) -- (-0.3,-0.8);
\node[anchor=north, text=blue!60, font=\footnotesize\bfseries]
    at (-2.0,-3.0) {\textbf{地球红外}};
\node[anchor=south, text=blue!60, font=\scriptsize]
    at (-1.2,-2.3) {$\sim$240\,W/m\textsuperscript{2}};

% ===== 地球反照（下方浅蓝，从地球指向卫星） =====
\draw[->, line width=1.2pt, color=cyan!60]
    (1.2,-3.0) -- (0.3,-0.8);
\node[anchor=north, text=cyan!60, font=\footnotesize\bfseries]
    at (2.0,-3.0) {\textbf{地球反照}};
\node[anchor=south, text=cyan!60, font=\scriptsize]
    at (1.2,-2.3) {$\sim$480\,W/m\textsuperscript{2}};

% ===== 辐射散热至深空（左右灰色箭头） =====
\draw[->, line width=1.5pt, color=graybase]
    (-1.5,0) -- (-3.5,0);
\node[anchor=east, text=graybase, font=\footnotesize\bfseries]
    at (-3.6,0.3) {\textbf{辐射散热}};
\node[anchor=east, text=graybase, font=\scriptsize]
    at (-3.6,-0.25) {1400\,W/m\textsuperscript{2}};

\draw[->, line width=1.5pt, color=graybase]
    (1.5,0) -- (3.5,0);
\node[anchor=west, text=graybase, font=\footnotesize\bfseries]
    at (3.6,0.3) {\textbf{辐射散热}};
\node[anchor=west, text=graybase, font=\scriptsize]
    at (3.6,-0.25) {1400\,W/m\textsuperscript{2}};

% ===== 深空标注 =====
\node[graybase, font=\scriptsize\itshape]
    at (3.8,2.5) {深空背景};
\node[graybase, font=\scriptsize\itshape]
    at (3.8,2.1) {2.7\,K};

% ===== 地球示意 =====
\filldraw[fill=blue!20, draw=blue!40, line width=0.8pt]
    (-3.0,-4.0) arc (180:360:3.0 and 0.8) -- (3.0,-4.0) -- cycle;
\node[blue!60, font=\scriptsize\bfseries] at (0,-4.0) {\textbf{地球}};

\end{tikzpicture}

\caption{LEO轨道热环境分析：太阳直射、地球红外辐射与反照}
\label{fig:leo_thermal}
\end{figure}

LEO轨道卫星面临三个主要热源：(1) 太阳直射（1361\,\unitWpm{}，定向后可降至约5\%吸收）；(2) 地球红外辐射（约240\,\unitWpm{}）；(3) 地球反照（约480\,\unitWpm{}）。三源合计即使经优化，仍对散热板形成约160--240\,\unitWpm{}的"寄生热流"。这意味着轨道散热不是"往冷空间一丢了之"，而是需要精确的热设计——散热面必须定向避开地球，太阳入射需通过刀口式（knife-edge）散热板最小化吸收面积。深空2.7\,K冷背景确实存在，但散热效率受限于中间介质——散热板的发射率、导热路径的热阻、以及热源（GPU）与辐射面之间的温度梯度。

\subsection{PUE视角：轨道与地面的能效对比}

传统数据中心用PUE（Power Usage Effectiveness）衡量能效。轨道数据中心PUE的等效估算见表~\ref{tab:pue}。

\begin{table}[H]
\centering
\footnotesize
\caption{轨道与地面数据中心PUE等效对比}
\label{tab:pue}
\begin{tabular}{lcc}
\toprule
\textbf{指标} & \textbf{轨道数据中心} & \textbf{地面数据中心（2026）} \\
\midrule
散热方式 & 被动辐射（零能耗） & 制冷机+冷塔（机械） \\
PUE（散热部分） & $\sim$1.00 & $\sim$1.05--1.10 \\
\addlinespace
供电方式 & 光伏+储能 & 双路电网 \\
PUE（供电部分） & $\sim$1.30--1.50 & $\sim$1.02--1.03 \\
\addlinespace
\textbf{综合等效PUE} & $\sim$\textbf{1.35} & $\sim$\textbf{1.10（行业平均）} \\
\bottomrule
\end{tabular}
\end{table}

轨道数据中心在散热端确实占优——被动辐射冷却不消耗额外能量。但供电端的光伏转换损失（$\sim$30\%）与储能循环损耗使得综合等效PUE反高于地面。这一逆直觉的结论是后续商业分析中轨道TCO偏高的物理根源之一。

\subsection{科学层小结}

物理学层面，太空数据中心散热方案在数学上可行，AI1卫星的工作点经斯特藩-玻尔兹曼方程验证成立。但当规模从单星推至GW集群时，散热面积需求呈线性增长（$A = P/q$），71万平米的辐射面构型超出了当前在轨大型结构的工程经验。PUE等效对比进一步揭示：轨道在散热端省去了机械制冷能耗，但供电端的光伏损失抵消了这一优势。"物理上可以，但余量有限"——这是科学层对马斯克散热主张的总体评价。
